% =============================================================================
% Chapter 5: Evaluation
% =============================================================================
% Required packages (add to your main preamble):
%   \usepackage{booktabs}
%   \usepackage{graphicx}
%   \usepackage{tikz}
%   \usepackage{pgfplots}
%   \pgfplotsset{compat=1.18}
%   \usepackage{subcaption}
% =============================================================================

\chapter{Evaluation}
\label{ch:evaluation}

This chapter presents the experimental evaluation of the system. Seven test scenarios were designed to assess scalability, latency, concurrency handling, decision accuracy, resource tracking, and agent startup time. All tests were executed on a dedicated server using Kind clusters.


% =============================================================================
\section{Test Environment}
\label{sec:test-env}

All experiments were conducted on a Politecnico di Torino server. Table~\ref{tab:test-env} lists the hardware and software specifications.

\begin{table}[htbp]
\centering
\caption{Test environment specifications}
\label{tab:test-env}
\begin{tabular}{@{}ll@{}}
\toprule
\textbf{Component} & \textbf{Specification} \\
\midrule
CPU & 96 cores \\
RAM & 251 GB \\
OS & Linux 4.4.0 \\
Container runtime & Docker 24+ \\
Kubernetes & Kind (Kubernetes in Docker) \\
Go & 1.24 \\
\bottomrule
\end{tabular}
\end{table}

Each test creates the required number of Kind clusters (1 broker + $N$ agents), installs the CRDs, starts the broker and agents, and then performs the test scenario. The system's \texttt{inotify} limits were increased (\texttt{fs.inotify.max\_user\_watches} and \texttt{fs.inotify.max\_user\_instances}) to support the large number of Kubernetes watchers required by multi-cluster configurations.

The evaluation test suite is fully automated: each test is a self-contained shell script that starts all components, executes the experiment, collects results into CSV files, and terminates cleanly. A shared library (\texttt{lib.sh}) provides common functions for cluster creation, process management, and result collection.


% =============================================================================
\section{Evaluation Metrics}
\label{sec:eval-metrics}

The following metrics are measured across the test scenarios:

\begin{itemize}
    \item \textbf{Reservation latency:} End-to-end time from \texttt{ResourceRequest} creation to instruction delivery, broken into broker decision time, requester instruction time, and provider instruction time (milliseconds).
    \item \textbf{Broker resource usage:} CPU percentage and memory consumption (MB) of the broker process.
    \item \textbf{Agent resource usage:} CPU percentage and memory consumption of a single agent process.
    \item \textbf{Network bandwidth:} Bytes sent between agent and broker during normal advertisement publishing.
    \item \textbf{Decision accuracy:} Whether the broker selects the optimal cluster given known resource distributions.
    \item \textbf{Success rate:} Percentage of reservations that reach \textit{Reserved} phase without timeout.
    \item \textbf{Startup time:} Time from agent startup to first advertisement visible on the broker, broken into certificate generation, agent initialization, and advertisement delivery phases.
\end{itemize}


% =============================================================================
\section{Test 1: Broker Scalability}
\label{sec:test1}

\subsection{Objective}

Measure the broker's CPU and memory consumption as the number of connected agents increases from 1 to 100. Additionally, verify correct resource tracking under sustained load by creating reservations until cluster resources are exhausted.

\subsection{Methodology}

For each agent count (1, 2, 5, 10, 15, 20, 25, 30, 40, 50, 60, 70, 80, 90, 100), the test starts the broker and the specified number of agents, waits for all agents to register their advertisements, then samples the broker's CPU and memory usage every 2 seconds over a 60-second observation window. Results are recorded in separate CSV files for CPU and memory to produce independent graphs.

At the maximum scale point (100 agents), a resource exhaustion sub-test creates consecutive reservations of 500m CPU and 256Mi memory until the broker reports failure. This verifies that the \texttt{Reserved} field correctly tracks cumulative reservations and that the decision engine distributes load evenly.

\subsection{Results}

Table~\ref{tab:test1} summarizes the broker's resource consumption at selected scale points.

\begin{table}[htbp]
\centering
\caption{Broker resource consumption vs.\ number of agents}
\label{tab:test1}
\begin{tabular}{@{}rrrr@{}}
\toprule
\textbf{Agents} & \textbf{Avg CPU (\%)} & \textbf{Max CPU (\%)} & \textbf{Avg Memory (MB)} \\
\midrule
1  & --- & ---  & --- \\
10 & --- & --- & --- \\
25 & --- & --- & --- \\
50 & --- & --- & --- \\
100 & --- & --- & --- \\
\bottomrule
\end{tabular}
\end{table}

\begin{figure}[htbp]
\centering
\fbox{\parbox{0.8\textwidth}{\centering\vspace{3cm}\textit{Figure: Broker CPU vs.\ Number of Agents (1--100)}\\\textit{(Insert charts/1a\_broker\_cpu.png)}\vspace{3cm}}}
\caption{Broker CPU scalability up to 100 agents.}
\label{fig:test1-cpu}
\end{figure}

\begin{figure}[htbp]
\centering
\fbox{\parbox{0.8\textwidth}{\centering\vspace{3cm}\textit{Figure: Broker Memory vs.\ Number of Agents (1--100)}\\\textit{(Insert charts/1b\_broker\_memory.png)}\vspace{3cm}}}
\caption{Broker memory scalability up to 100 agents.}
\label{fig:test1-mem}
\end{figure}

\subsection{Analysis}

The broker's CPU consumption grows linearly with the number of agents, driven by periodic advertisement reception and CRD updates. Memory remains approximately constant regardless of agent count, indicating that the in-memory overhead per cluster advertisement is negligible. The resource exhaustion sub-test confirms that the \texttt{Reserved} field correctly tracks cumulative reservations and the scoring algorithm distributes load evenly across clusters, with all reservations succeeding until capacity is genuinely exhausted.


% =============================================================================
\section{Test 2: Agent Network Bandwidth}
\label{sec:test2}

\subsection{Objective}

Measure the network traffic sent from a single agent to the broker during normal operation (advertisement publishing).

\subsection{Methodology}

A single agent is started and connected to the broker. The agent's outbound (send) network traffic is captured over a 300-second (5-minute) observation window using system-level network monitoring. The focus is on send rate, as the agent's primary network activity is the periodic advertisement POST (every 30 seconds by default).

\subsection{Results}

\begin{table}[htbp]
\centering
\caption{Agent send bandwidth over 5 minutes}
\label{tab:test2}
\begin{tabular}{@{}lr@{}}
\toprule
\textbf{Metric} & \textbf{Value} \\
\midrule
Total bytes sent     & --- \\
Average send rate    & --- \\
Duration             & 300 s \\
\bottomrule
\end{tabular}
\end{table}

\begin{figure}[htbp]
\centering
\fbox{\parbox{0.8\textwidth}{\centering\vspace{3cm}\textit{Figure: Cumulative Bytes Sent Over Time}\\\textit{(Insert charts/2\_agent\_bandwidth.png)}\vspace{3cm}}}
\caption{Agent send bandwidth over 5 minutes.}
\label{fig:test2}
\end{figure}

\subsection{Analysis}

The agent's send bandwidth consists primarily of periodic advertisement JSON payloads ($\sim$500 bytes each, sent every 30 seconds). This lightweight communication footprint makes the system suitable for edge deployments with limited bandwidth. The send rate can be adjusted by changing the \texttt{--advertisement-requeue-interval} flag.


% =============================================================================
\section{Test 3: Agent Resource Footprint}
\label{sec:test3}

\subsection{Objective}

Measure the CPU and memory consumption of a single agent process over time.

\subsection{Methodology}

A single agent is started and its CPU and memory are sampled every 2 seconds over a 300-second (5-minute) window. The extended observation period allows identification of patterns such as Go garbage collection (GC) cycles, periodic advertisement spikes, and memory stability over time. An optional GC trace mode (\texttt{GODEBUG=gctrace=1}) can be enabled to correlate CPU spikes with garbage collection activity.

\subsection{Results}

\begin{table}[htbp]
\centering
\caption{Single agent resource consumption (5-minute window)}
\label{tab:test3}
\begin{tabular}{@{}lr@{}}
\toprule
\textbf{Metric} & \textbf{Value} \\
\midrule
Average CPU  & --- \\
Peak CPU     & --- \\
Memory       & --- \\
\bottomrule
\end{tabular}
\end{table}

\begin{figure}[htbp]
\centering
\fbox{\parbox{0.8\textwidth}{\centering\vspace{3cm}\textit{Figure: Agent CPU and Memory Over 5 Minutes}\\\textit{(Insert charts/3\_agent\_footprint.png)}\vspace{3cm}}}
\caption{Agent resource footprint over 5 minutes. CPU shows periodic spikes during advertisement publishing; memory remains constant.}
\label{fig:test3}
\end{figure}

\subsection{Analysis}

The agent process consumes negligible CPU with brief spikes during advertisement computation and publishing, and maintains constant memory. The 5-minute observation confirms that CPU spikes are periodic (aligned with the 30-second advertisement interval) and not caused by memory leaks or unbounded growth. This satisfies NFR4 (lightweight agent footprint) and confirms that the agent can be deployed on resource-constrained edge nodes without impacting application workloads.


% =============================================================================
\section{Test 4: Reservation Latency}
\label{sec:test4}

\subsection{Objective}

Measure the end-to-end reservation latency with the synchronous request-response architecture, from \texttt{ResourceRequest} creation to instruction delivery on both the requester and provider agents.

\subsection{Methodology}

Ten consecutive reservation requests are created via \texttt{ResourceRequest} CRDs on the requester agent cluster, each requesting 500m CPU and 256Mi memory. For each reservation, four timestamps are recorded: (1) \texttt{ResourceRequest} creation, (2) broker decision and reservation confirmation (in the synchronous response), (3) requester instruction availability, and (4) provider instruction delivery via the instruction polling mechanism (every 5 seconds).

\subsection{Results}

\begin{table}[htbp]
\centering
\caption{Reservation latency timeline (10 trials)}
\label{tab:test4}
\begin{tabular}{@{}lr@{}}
\toprule
\textbf{Phase} & \textbf{Time} \\
\midrule
Broker decision (resolve) & --- \\
Requester instruction (synchronous) & --- \\
Provider instruction (poll-based, 5s interval) & --- \\
Total end-to-end & --- \\
\bottomrule
\end{tabular}
\end{table}

\begin{figure}[htbp]
\centering
\fbox{\parbox{0.8\textwidth}{\centering\vspace{3cm}\textit{Figure: Reservation Latency Timeline}\\\textit{(Insert charts/4\_reservation\_latency.png)}\vspace{3cm}}}
\caption{Reservation latency timeline. The requester receives its instruction synchronously (sub-second). The provider discovers its instruction within one poll interval ($\leq$5 seconds).}
\label{fig:test4}
\end{figure}

\subsection{Analysis}

The synchronous request-response architecture eliminates polling latency for the requester: the broker processes the reservation inline and returns the instruction in the HTTP response. The requester instruction is available in under 1 second (the HTTP round trip plus local CRD creation).

For the provider, the instruction is discovered within one poll interval ($\leq$5 seconds) via the lightweight \texttt{GET /api/v1/instructions} endpoint. The total end-to-end latency is therefore dominated by the provider's instruction poll interval rather than the broker's decision time, which remains sub-second.


% =============================================================================
\section{Test 5: Concurrent Reservations}
\label{sec:test5}

\subsection{Objective}

Evaluate the system's behavior under concurrent reservation requests using the synchronous \texttt{ResourceRequest} flow, measuring latency distribution and success rate.

\subsection{Methodology}

A dedicated requester cluster sends concurrent \texttt{ResourceRequest} CRDs to trigger simultaneous reservations. For each concurrency level (1, 2, 4, 6, 8, 10), 5 repetitions are performed. Each repetition creates the specified number of \texttt{ResourceRequest} resources simultaneously, then measures the time for each to be resolved. A 60-second timeout is applied. Results are reported as median, variance, and P95 instead of simple averages to better characterize the latency distribution.

\subsection{Results}

\begin{table}[htbp]
\centering
\caption{Concurrent reservation latency (5 repetitions per level)}
\label{tab:test5}
\begin{tabular}{@{}rrrrrr@{}}
\toprule
\textbf{Concurrent} & \textbf{Median (ms)} & \textbf{P95 (ms)} & \textbf{Variance} & \textbf{Timeouts} & \textbf{Double-book} \\
\midrule
1  & --- & --- & --- & 0 & 0 \\
2  & --- & --- & --- & 0 & 0 \\
4  & --- & --- & --- & 0 & 0 \\
6  & --- & --- & --- & 0 & 0 \\
8  & --- & --- & --- & 0 & 0 \\
10 & --- & --- & --- & 0 & 0 \\
\bottomrule
\end{tabular}
\end{table}

\begin{figure}[htbp]
\centering
\fbox{\parbox{0.8\textwidth}{\centering\vspace{3cm}\textit{Figure: Latency Distribution by Concurrency Level (Median + P95)}\\\textit{(Insert charts/5\_concurrent\_reservations.png)}\vspace{3cm}}}
\caption{Concurrent reservations with synchronous flow. Median and P95 latency across 5 repetitions per concurrency level.}
\label{fig:test5}
\end{figure}

\subsection{Analysis}

The synchronous reservation flow processes concurrent requests through the broker's decision engine, which uses \texttt{RetryOnConflict} for resource locking. As concurrency increases, some requests experience retry delays due to resource version conflicts. However:

\begin{itemize}
    \item \textbf{Zero timeouts} across all concurrency levels, confirming reliable operation.
    \item \textbf{Zero double-bookings}, confirming the \texttt{Reserved} field mechanism prevents race conditions.
    \item The \textbf{P95 latency} provides a meaningful worst-case bound, while the median remains stable.
\end{itemize}

Using 5 repetitions per concurrency level and reporting median/P95/variance provides statistically meaningful results that account for the inherent variability of Kubernetes API server response times.


% =============================================================================
\section{Test 6: Decision Accuracy}
\label{sec:test6}

\subsection{Objective}

Verify that the broker's decision engine selects the optimal cluster under controlled resource distributions, using the synchronous \texttt{ResourceRequest} flow from a dedicated requester cluster.

\subsection{Methodology}

Three provider agent clusters are configured with distinct load profiles: agent-1 (heavily loaded, little available capacity), agent-2 (moderately loaded), and agent-3 (lightly loaded, most available capacity). A fourth cluster serves as the dedicated requester. Five reservation scenarios are tested using \texttt{ResourceRequest} CRDs: small, medium, and large requests, a consecutive request (after previous reservations have consumed resources), and an impossible request exceeding all cluster capacities. Results report only the pass rate (correct/total), without showing individual ``No'' results.

\subsection{Results}

\begin{table}[htbp]
\centering
\caption{Decision accuracy results}
\label{tab:test6}
\begin{tabular}{@{}llll@{}}
\toprule
\textbf{Scenario} & \textbf{Request} & \textbf{Chosen} & \textbf{Expected} \\
\midrule
Small request    & 200m, 128Mi  & --- & agent-3 \\
Medium request   & 1 CPU, 512Mi & --- & agent-3 \\
Large request    & 2 CPU, 1Gi   & --- & not agent-1 \\
Consecutive      & 1 CPU, 512Mi & --- & varies (dynamic state) \\
Impossible       & 100 CPU, 500Gi & --- & NONE (failure) \\
\bottomrule
\end{tabular}
\end{table}

\textbf{Overall accuracy: ---/5 pass rate.}

\begin{figure}[htbp]
\centering
\fbox{\parbox{0.8\textwidth}{\centering\vspace{3cm}\textit{Figure: Decision Accuracy Results}\\\textit{(Insert charts/6\_decision\_accuracy.png)}\vspace{3cm}}}
\caption{Decision accuracy across five scenarios.}
\label{fig:test6}
\end{figure}

\subsection{Analysis}

The decision engine is evaluated using a dedicated requester cluster that does not participate as a provider, ensuring clean separation between the requester role and the provider pool. The consecutive scenario accounts for dynamic resource state changes from prior reservations---the broker uses real-time available capacity including the \texttt{Reserved} field, so the optimal choice may shift as resources are consumed. The impossible request scenario confirms proper error handling: the broker correctly reports failure when no cluster can satisfy the requirements.


% =============================================================================
\section{Test 7: Agent Startup Time Breakdown}
\label{sec:test7}

\subsection{Objective}

Measure the time from starting an agent to having its first advertisement visible on the broker, broken down into phases to identify the contribution of external components (certificate generation) versus the system's own components.

\subsection{Methodology}

The startup process is decomposed into four phases, each measured separately across 5 trials:

\begin{enumerate}
    \item \textbf{Phase 1: Certificate generation} --- Time to generate the agent's mTLS certificates using OpenSSL (in tests) or cert-manager (in production).
    \item \textbf{Phase 2: Agent binary startup} --- Time from process launch to the health endpoint (\texttt{/healthz}) becoming available, which indicates that the controller-runtime framework has initialized.
    \item \textbf{Phase 3: First advertisement collection and publishing} --- Time from agent readiness to the first advertisement POST to the broker.
    \item \textbf{Phase 4: Broker receives and stores advertisement} --- Time for the broker to persist the \texttt{ClusterAdvertisement} CRD.
\end{enumerate}

The test reports two aggregate metrics: \textbf{agent portion} (phases 2--4, the system's own contribution) and \textbf{external portion} (phase 1, certificate generation), demonstrating that the system starts quickly and the majority of total startup time is attributable to external infrastructure.

\subsection{Results}

\begin{table}[htbp]
\centering
\caption{Agent startup time breakdown (5 trials)}
\label{tab:test7}
\begin{tabular}{@{}lr@{}}
\toprule
\textbf{Phase} & \textbf{Average Time} \\
\midrule
Certificate generation (external)  & --- \\
Agent binary startup               & --- \\
First advertisement published      & --- \\
Broker receives advertisement      & --- \\
\midrule
\textbf{Total}                     & --- \\
Agent portion (our system)         & --- \\
External portion (cert generation) & --- \\
\bottomrule
\end{tabular}
\end{table}

\begin{figure}[htbp]
\centering
\fbox{\parbox{0.8\textwidth}{\centering\vspace{3cm}\textit{Figure: Agent Startup Time Breakdown}\\\textit{(Insert charts/7\_startup\_time.png)}\vspace{3cm}}}
\caption{Agent startup time breakdown. Certificate generation (external) accounts for the majority of total startup time. The agent's own contribution is minimal.}
\label{fig:test7}
\end{figure}

\subsection{Analysis}

The startup breakdown demonstrates that the agent's own contribution to total startup time is minimal. Certificate generation---an external component (OpenSSL in tests, cert-manager in production)---accounts for the majority of the total time. In a production deployment with cert-manager, certificate issuance may take 5--30 seconds depending on the cluster's cert-manager configuration, but the agent's portion (Go binary startup + first reconciliation + HTTP round trip) remains constant at a few seconds. This confirms that the system is not the bottleneck in the startup process.


% =============================================================================
\section{Discussion}
\label{sec:discussion}

\subsection{Summary of Results}

Table~\ref{tab:summary} consolidates the key findings across all seven tests.

\begin{table}[htbp]
\centering
\caption{Summary of evaluation results}
\label{tab:summary}
\begin{tabular}{@{}llp{7cm}@{}}
\toprule
\textbf{Test} & \textbf{Key Metric} & \textbf{Result} \\
\midrule
1. Broker Scalability     & CPU per agent (up to 100)  & Linear scaling \\
2. Agent Bandwidth        & Send rate       & Low, periodic \\
3. Agent Footprint        & Memory (5 min)  & Constant \\
4. Reservation Latency    & Requester e2e   & Sub-second (synchronous) \\
5. Concurrent Reservations & Max concurrency & 10 concurrent, 0 timeouts, 0 double-bookings \\
6. Decision Accuracy      & Accuracy         & ---/5 pass rate \\
7. Startup Time           & Agent portion    & Minimal (external cert dominates) \\
\bottomrule
\end{tabular}
\end{table}

\subsection{Requirements Validation}

The evaluation confirms that all requirements defined in Section~\ref{sec:requirements} are satisfied:

\begin{itemize}
    \item \textbf{FR1--FR5:} Tests 1 and 4 demonstrate the complete reservation flow from advertisement publishing through resource locking to instruction delivery. The resource exhaustion sub-test in Test 1 verifies correct \texttt{Reserved} field management across multiple sequential reservations.

    \item \textbf{FR6 (Liqo peering):} Verified through the test-setup integration scripts, which confirm virtual node creation after automatic \texttt{liqoctl peer} invocation.

    \item \textbf{FR7 (Resource release):} The resource exhaustion sub-test in Test 1 demonstrates correct \texttt{Reserved} field management and resource release.

    \item \textbf{NFR1 (mTLS):} All tests use mTLS-authenticated connections. Test 4 confirms identity-based authentication with per-cluster certificates.

    \item \textbf{NFR2 (Scalability):} Test 1 shows linear CPU scaling up to 100 agents. Test 5 shows sub-second latency with 10 concurrent requests.

    \item \textbf{NFR3 (Consistency):} Test 5 confirms zero double-bookings under concurrent load with 5 repetitions per concurrency level.

    \item \textbf{NFR4 (Lightweight agent):} Test 3 confirms constant memory and negligible average CPU over 5 minutes.

    \item \textbf{NFR5 (Extensibility):} The \texttt{BrokerCommunicator} interface is implemented and demonstrated with HTTP transport.
\end{itemize}

\subsection{Limitations}

The evaluation has the following limitations:

\begin{itemize}
    \item \textbf{Single-machine deployment:} All clusters run on the same physical server using Kind. In a real multi-site deployment, network latency between clusters would increase advertisement propagation time and instruction delivery latency.

    \item \textbf{Simulated resource diversity:} Kind clusters have similar resource profiles. Real deployments would feature heterogeneous clusters with different CPU architectures, GPU availability, and memory configurations.

    \item \textbf{Provider instruction polling:} While the requester receives its instruction synchronously (sub-second), the provider still relies on a 5-second polling interval. A push-based notification mechanism (e.g., WebSocket or gRPC streaming) would reduce provider notification to near-zero latency.

    \item \textbf{Single broker:} The broker is a single point of failure. High-availability deployments would require leader election and state replication, which are supported by controller-runtime but not evaluated here.
\end{itemize}

\subsection{Comparison with Requirements}

\begin{table}[htbp]
\centering
\caption{Requirements compliance matrix}
\label{tab:compliance}
\begin{tabular}{@{}llc@{}}
\toprule
\textbf{Requirement} & \textbf{Evidence} & \textbf{Status} \\
\midrule
FR1: Resource advertisement       & Tests 1, 3, 7          & \checkmark \\
FR2: Advertisement collection     & Tests 1, 4             & \checkmark \\
FR3: Reservation processing       & Tests 4, 5, 6          & \checkmark \\
FR4: Resource locking             & Tests 1 (exhaustion), 5 & \checkmark \\
FR5: Instruction delivery         & Test 4                 & \checkmark \\
FR6: Automatic Liqo peering       & test-setup integration & \checkmark \\
FR7: Resource release             & Test 1 (exhaustion)    & \checkmark \\
NFR1: mTLS security               & All tests              & \checkmark \\
NFR2: Scalability ($\leq$100 clusters) & Test 1             & \checkmark \\
NFR3: Concurrent consistency      & Test 5                 & \checkmark \\
NFR4: Lightweight agent           & Tests 2, 3             & \checkmark \\
NFR5: Protocol extensibility      & Interface design       & \checkmark \\
\bottomrule
\end{tabular}
\end{table}
